%%
%% This is file `TikZ_Studio_en.tex',
%% English user guide for TikZ Studio.
%%
%% Copyright (C) 2026 Shijin Chen
%%
%% Build: pdflatex TikZ_Studio_en && pdflatex TikZ_Studio_en
%%
\documentclass[a4paper]{ltxdoc}
\usepackage[T1]{fontenc}
\usepackage{lmodern}
\usepackage{booktabs}
\usepackage{array}
\usepackage{longtable}
\usepackage{enumitem}
\usepackage{xcolor}
\usepackage{hyperref}
\usepackage{url}

\hypersetup{
  colorlinks=true,
  linkcolor=blue!50!black,
  urlcolor=blue!50!black,
  pdftitle={TikZ Studio --- User Guide (English)},
  pdfauthor={Shijin Chen},
  pdfsubject={Local TikZ / circuitikz editor with interactive preview},
}

\newcommand{\fileversion}{2026-11}
\newcommand{\filedate}{2026-11-21}
\newcommand{\appname}{\textsf{TikZ Studio}}
\newcommand{\ui}[1]{\textbf{#1}}
\newcommand{\key}[1]{\texttt{#1}}

\title{%
  \appname\ --- User Guide%
  \thanks{This file describes version~\fileversion{} (\filedate).}%
}
\author{%
  Shijin Chen\\
  \texttt{sjchen1998@gmail.com}%
}
\date{Printed \today\\[0.5em] Guide date: \filedate{} (v\fileversion)}

\begin{document}
\maketitle

\begin{abstract}
\appname\ is a local desktop / browser application for editing
\textbf{TikZ} and \textbf{circuitikz}.  Source is on the left; XeLaTeX or
pdfLaTeX draws the preview on the right.  You place components by
clicking the canvas, drag red handles to rewrite coordinates, box-select
and transform groups, edit dimensions after the fact, and import DXF
into geometry or \textsf{satopdf} structural figures.  This guide focuses
on \emph{how to edit}.
\end{abstract}

\tableofcontents
\clearpage

%==============================================================================
\section{Introduction}
%==============================================================================

\appname\ targets lecture notes, reports, and a CAE$\rightarrow$PDF
pipeline on one machine.  Licence: GPL-3.0 (\texttt{LICENSE}).
Copyright \textcopyright\ 2026 Shijin Chen.

Related tools in this repository: \path{lisp/LoadXData.lsp} (AutoCAD
annotation), \path{texmf/tex/latex/local/satopdf.sty} (structural TikZ).

%==============================================================================
\section{Install and start}
%==============================================================================

\begin{itemize}[nosep]
  \item Python~3; \texttt{pip install -r requirements.txt}
        (\texttt{fastapi}, \texttt{uvicorn}, \texttt{python-multipart},
        \texttt{pymupdf}, \texttt{pywebview}, \texttt{ezdxf});
  \item TeX~Live with XeLaTeX, pdfLaTeX, \textsf{ctex}, \textsf{circuitikz},
        \textsf{tikz-3dplot}, \textsf{pgfplots};
  \item In-tree \textsf{satopdf} under \path{texmf/tex/latex/local/}.
\end{itemize}

\begin{verbatim}
python -m tikzstudio          # desktop window (recommended)
python -m tikzstudio.server   # browser only -> http://127.0.0.1:8765
\end{verbatim}
Or double-click \path{start_tikz_studio.bat} (always starts from its own
folder).

%==============================================================================
\section{Window layout}
%==============================================================================

\begin{itemize}[nosep]
  \item \textbf{Left --- toolbar:} Compile, Undo, Open \texttt{.tex},
        Import DXF, Save, Template menu, compiler (XeLaTeX/pdfLaTeX),
        language (EN/ZH), status text.
  \item \textbf{Left --- component bar:} buttons for the current panel,
        plus a \ui{+} library button.
  \item \textbf{Left --- editor:} TeX source; selected preview points
        highlight the matching line.
  \item \textbf{Right --- preview bar:} zoom, point pick on/off, show
        red/blue/purple, Link, Move-all, snap grid, ortho-snap, export
        PDF/PNG/SVG.
  \item \textbf{Right --- edit bar:} Select-all, placement prompts,
        dimension style controls, X/Y (or $\Delta$X/$\Delta$Y) inputs,
        Apply, Delete, Scale/Rotate/Transform.
  \item \textbf{Right --- canvas:} SVG preview with draggable overlays.
\end{itemize}

%==============================================================================
\section{Basic editing loop}
%==============================================================================

\begin{enumerate}[nosep]
  \item Pick a \ui{template} (or keep the current figure).
  \item Click a component on the bar $\rightarrow$ placement mode starts
        (banner on the edit bar).
  \item Click the required points on the preview (or fill form fields for
        0-point tools).
  \item Press \ui{Confirm}, press Enter in a field, or
        \textbf{right-click} the canvas to finish.
  \item \ui{Compile} (or wait for auto-compile after a successful place)
        to refresh handles.
  \item Drag red points, or select several and transform; export when done.
\end{enumerate}

\noindent\textbf{Cancel placement:} \key{Esc}, or \ui{Cancel}.\\
\textbf{Compile from the editor:} \key{Ctrl}+\key{Enter}.\\
\textbf{Undo:} toolbar \ui{Undo} (restores previous buffer snapshot).\\
\textbf{Failed compile:} the editor reverts to the last source that
compiled cleanly.

%==============================================================================
\section{Templates and picture panels}
%==============================================================================

The component bar follows the \emph{picture kind} inferred from the
source:

\begin{center}
{\small
\begin{tabular}{@{}llp{0.40\linewidth}@{}}
\toprule
Kind & Detected by & Core tools \\
\midrule
Circuit & \texttt{circuitikz} & Text, wire, R/C/L, V/I, switch, dims \\
Blocks & named \texttt{draw} nodes & Box, round, ellipse, cylinder,
  trap, decision, arrows, fill \\
Plot & \texttt{axis}/\texttt{addplot} & Function, marks, label, dims \\
Plane & plane-style \texttt{tikzpicture} & Point, segment, function, dims \\
3D & \texttt{tikz-3dplot} markers & 3D segment/arrow/circle/arc, dims \\
\bottomrule
\end{tabular}}
\end{center}

Built-in templates: RC filter, Buck, op-amp, block diagram, \textbf{2D SA},
\textbf{3D SA}, plot, 2D axes, CMS~3D, Gaussian, frequency response.
2D/3D SA use \textsf{satopdf}.

Changing the template \emph{replaces} the editor buffer.  Save first if
you need the old figure.

%==============================================================================
\section{How to place components}
\label{sec:place}
%==============================================================================

\subsection{General rule}
Each tool declares how many canvas clicks it needs ($n=0$, $1$, $2$,
or more for polylines / dimensions):

\begin{itemize}[nosep]
  \item $n=0$ --- fill the form on the edit bar, then \ui{Confirm}
        (examples: plot function, plane axes ranges, 3D circle fields).
  \item $n=1$ or $2$ --- click that many points; the live title shows
        ``click $k/n$''.  After the last point you can still edit fields
        and press Enter / \ui{Confirm}.
  \item Snap: clicks prefer nearby named ports / coordinates when
        available; otherwise numeric TikZ units under the cursor.
\end{itemize}

Inserted code is appended just before
\texttt{\textbackslash end\{tikzpicture\}},
\texttt{\textbackslash end\{circuitikz\}}, or
\texttt{\textbackslash end\{axis\}} (for plot tools).

\subsection{Circuit panel}
Click two endpoints for two-terminal parts (R, C, L, V, I, switch,
diode, meter, \ldots).  One-terminal devices (ground, MOS, op-amp, \ldots)
need a single placement point and optional rotation / label.  Labels go
into the text field before Confirm.

\subsection{Block diagram}
\begin{itemize}[nosep]
  \item Shapes with text require an \textbf{ID} (letter + alphanumerics).
        A duplicate shows ``ID already exists''.
  \item After compile, ports appear on named nodes --- use them when
        placing arrows so the source stores \texttt{(id.east)} style
        anchors instead of raw numbers when possible.
  \item \ui{Polyline arrow}: click start, then corners; finish with
        right-click or \ui{Confirm} (minimum two points).  With
        \ui{Ortho snap} on, nearly horizontal/vertical segments lock.
\end{itemize}

\subsection{Plot / plane / 3D field tools}
Enter expressions, domains, or 3D centre/plane strings in the form.
For 3D circle/arc, \texttt{plane} may be \texttt{xy}/\texttt{xz}/\texttt{yz},
a normal, or two directions (first direction $=0^\circ$).  Centre may be
\texttt{(0,0,1.55)} or a name such as \texttt{(O)}.

\subsection{Shared geometry (most panels)}
Available from the bar or via \ui{+} library: line, arrow, dashed,
rectangle (two corners), fill, grid, circle, arc, ellipse path,
coordinate, \ldots

\paragraph{Circle.}
Click centre, then a second point (or hover a tangent line / curve /
through-point); radius updates live.  Second click confirms.

\paragraph{Arc.}
Centre $\rightarrow$ start $\rightarrow$ end (angles computed).

\paragraph{Ellipse path.}
Centre $\rightarrow$ horizontal-axis end $\rightarrow$ vertical-axis end.

\subsection{Dimensions}
\begin{itemize}[nosep]
  \item \ui{Dimension} --- click start, end, then label position.
        While placing (after two points):
        \key{H}/\key{V}/\key{R} = horizontal / vertical / aligned;
        \key{I}/\key{O} = arrows inside / outside;
        \key{1}/\key{2} = hide start / end extension.
  \item \ui{Radius} / \ui{Diameter} --- click an existing circle or arc,
        then the label position.
  \item \ui{Angle} --- click first line, second line, then label
        position; \key{I}/\key{O} for arrow style.
\end{itemize}
After placement, select a dimension handle to reopen the style bar
(horizontal / vertical / aligned, in/out, suppress sides, label).
Press Enter in the label field to apply.

\subsection{Custom library}
Open \ui{+}:
\begin{enumerate}[nosep]
  \item Add a built-in tool that is not pinned yet, or
  \item Define a custom button: name, click count (0/1/2), optional text
        field, and a TikZ template using
        \texttt{\{\{1\}\}}, \texttt{\{\{2\}\}},
        \texttt{\{\{text\}\}} / Chinese field name, and
        \texttt{\{\{angle\}\}}.
\end{enumerate}
Library layout is stored in the browser (\texttt{localStorage}).

%==============================================================================
\section{How to edit existing geometry}
\label{sec:edit}
%==============================================================================

\subsection{Handle colours}
\begin{center}
\begin{tabular}{@{}ll@{}}
\toprule
Colour & Meaning \\
\midrule
Red & Absolute numeric coordinate --- drag to rewrite the source \\
Blue & Named node / relative / port-like position \\
Purple & Ratio / interpolated point \\
\bottomrule
\end{tabular}
\end{center}

Turn \ui{Point pick} off to skip parsing handles (faster on huge figures).
Toggle red/blue/purple visibility independently.

\subsection{Drag one point}
\begin{enumerate}[nosep]
  \item Ensure \ui{Point pick} is on and red points are visible.
  \item Press and drag a red handle; release to write the new
        \texttt{(x,y)} into the source and recompile.
  \item \ui{Link} (default on): other occurrences of the same numeric pair
        move together.
  \item \ui{Move all}: while dragging one absolute point, every numeric
        pair on the same TeX statement (line through \texttt{;}) shifts
        by the same $\Delta$.
\end{enumerate}
Snap grid: free / $0.05$ / $0.1$ / $0.25$.

\subsection{Numeric fields (single selection)}
Click a handle.  The edit bar shows \ui{X}/\ui{Y} (TikZ units).  Edit and
press \ui{Apply} or Enter.  Special cases:
\begin{itemize}[nosep]
  \item Ellipse path --- also edit semi-axes and rotation.
  \item Packaged plane axes --- edit axis and tick ranges in the extra
        fields (handles themselves are not box-selected).
\end{itemize}

\subsection{Multi-select, box select, select all}
\begin{itemize}[nosep]
  \item \key{Ctrl}-click or \key{Shift}-click adds/removes handles.
  \item Drag on empty canvas to rubber-band select.
  \item \ui{Select all} or \key{Ctrl}+\key{A} (when focus is not in an
        input) selects all movable handles.
  \item With a multi-selection, X/Y become \textbf{$\Delta$X / $\Delta$Y};
        \ui{Apply} translates the group.
\end{itemize}

\noindent\textbf{Protected (never multi-selected):}
\texttt{planeaxes} and grid points; dimension / radius / diameter /
angle construction points.  Edit those via their dedicated fields or
style bar.

\subsection{Scale and rotate}
With a selection, enter \ui{Scale} and \ui{Rotate} (degrees), then
\ui{Transform}.
\begin{itemize}[nosep]
  \item After \ui{Select all}, scale applies to the whole picture:
        coordinates, circle/arc radii, \texttt{--++} lengths, and
        \texttt{\textbackslash tdplotsetcoord} vector lengths.
  \item A partial box only rewrites numeric coordinates inside the box.
\end{itemize}

\subsection{Delete}
With an active selection, \ui{Delete} removes the corresponding source
construct when supported (e.g.\ a placed statement / scope).  Prefer
Undo if unsure.

\subsection{Source editor}
\begin{itemize}[nosep]
  \item Edit TikZ by hand at any time; \ui{Compile} or
        \key{Ctrl}+\key{Enter} to refresh.
  \item \key{Tab} / \key{Shift}+\key{Tab} indent / dedent (width from
        \texttt{.editorconfig}).
  \item Selecting a preview point scrolls/highlights its source line.
\end{itemize}

%==============================================================================
\section{Preview navigation}
%==============================================================================

\begin{itemize}[nosep]
  \item Zoom slider $5\%$--$800\%$; mouse wheel zooms toward the cursor.
  \item Click empty preview space to clear the selection.
\end{itemize}

%==============================================================================
\section{Import DXF}
%==============================================================================

\ui{Import DXF} follows the \emph{active template}:

\begin{itemize}[nosep]
  \item \textbf{2D SA} --- \textsf{satopdf} with supports/loads from
        \texttt{LoadXData} XData;
  \item \textbf{3D SA} --- shells $\rightarrow$
        \texttt{\textbackslash stpDface}; face loads $\rightarrow$
        \texttt{\textbackslash stpDfaceload};
  \item \textbf{Other} --- geometry only; skip \texttt{CAE\_VIS\_*};
        scale large drawings to $\approx$16\,cm, origin at lower-left.
\end{itemize}
Limit 20\,MB.  For CAE data, open the SA template \emph{before} importing.

%==============================================================================
\section{Save and export}
%==============================================================================

\begin{itemize}[nosep]
  \item \ui{Save} --- desktop: Save~As \texttt{.tex}; browser: download.
  \item \ui{PDF} / \ui{PNG} / \ui{SVG} --- PNG at 600\,dpi.
\end{itemize}

%==============================================================================
\section{Keyboard and mouse cheatsheet}
%==============================================================================

\begin{center}
{\small
\begin{tabular}{@{}ll@{}}
\toprule
Input & Action \\
\midrule
\key{Ctrl}+\key{Enter} (editor) & Compile \\
\key{Esc} & Cancel placement / close library \\
Right-click (while placing) & Confirm place \\
Enter (place/edit fields) & Confirm / Apply \\
\key{Ctrl}/\key{Shift}-click handle & Multi-select \\
Drag empty canvas & Box select \\
\key{Ctrl}+\key{A} (outside inputs) & Select all handles \\
Wheel on preview & Zoom at cursor \\
\key{H}/\key{V}/\key{R} (dim place) & Horizontal / vertical / aligned \\
\key{I}/\key{O} (dim/angle place) & Arrows in / out \\
\key{1}/\key{2} (dim place) & Hide start / end side \\
\bottomrule
\end{tabular}}
\end{center}

%==============================================================================
\section{Recommended workflows}
%==============================================================================

\subsection{Circuit from scratch}
Template \ui{RC filter} (or empty circuitikz) $\rightarrow$ place wires and
R/C $\rightarrow$ drag red ends to tidy $\rightarrow$ add dimensions $\rightarrow$
export PDF.

\subsection{Block diagram}
Template \ui{Blocks} $\rightarrow$ place boxes with unique IDs $\rightarrow$
Compile once so ports appear $\rightarrow$ arrows / polyline arrows snapping
to ports $\rightarrow$ export.

\subsection{CAE figure}
AutoCAD + \texttt{LoadXData} $\rightarrow$ DXF $\rightarrow$ open
\ui{2D SA} or \ui{3D SA} $\rightarrow$ Import DXF $\rightarrow$ adjust with
handles / \textsf{satopdf} source $\rightarrow$ export.

%==============================================================================
\section{Troubleshooting}
%==============================================================================

\begin{itemize}[nosep]
  \item No handles --- enable \ui{Point pick} and red points; compile
        successfully first.
  \item Placement clicks ignored --- finish or \key{Esc} the current
        place mode; ensure compile produced a calibration (cal).
  \item Drag does nothing --- blue/purple points are not freely rewritten
        like red absolutes; edit names in source or drag a related red
        point.
  \item \textsf{satopdf} not found --- start from the project root so
        \texttt{TEXINPUTS} includes \path{texmf/tex/latex/local/}.
  \item Slow large figure --- disable point pick; hide unused colours.
\end{itemize}

\end{document}
